\documentclass[11pt]{scrartcl}
\usepackage[sexy]{evan}
\usepackage{answers}

\begin{document}
\title{Nepal PreTST 2026 Solutions}
\author{Prakrit Gajurel}
\date{28 June 2026}
\maketitle

\section{Problems}

\begin{problem}
    Let $a_1,a_2,\dots,a_{2n+2}$ be positive real numbers. Prove that they can be permuted into numbers $u_0,v_0,u_1,v_1,\dots,u_n,v_n$ such that 
    the polynomial
    $$P(x)=u_0-v_0x+u_1x^2-v_1x^3+\dots+u_nx^{2n}-v_nx^{2n+1}$$
    has at least one real root, and every real root of $P$ is greater than or equal to $1$.
\end{problem}

\begin{problem}
    Find all functions $f:\mathbb{Z}_{\ge 0}\to\mathbb{Z}_{\ge 0}$ such that for all 
    $m\in\mathbb{N}$,$$f(m-1)\mid m \qquad \text{ and } \qquad m+1 \mid T_{2026}(m),
    $$where $T_1(m)=f(m)$ and $T_{k+1}(m)=f\bigl(m+T_k(m)\bigr)$ for all $k\ge 1$.
\end{problem}

\begin{problem}
    Let $ABC$ be an acute scalene triangle with $\angle BAC=60^\circ$. Denote by $O$ and $H$ the circumcenter and orthocenter of
    $\triangle ABC$ respectively. Let $M$ be the midpoint of arc $BC$ on $(ABC)$ not containing $A$, and let ray $\overrightarrow{MH}$
    intersect $(ABC)$ a second time at $P$. Show that the lines $AP$, $OH$ and $BC$ pass through a common point.
\end{problem}

\begin{problem}
    Consider the standard triangular grid of unit equilateral triangles forming a larger equilateral triangle of side length $n$. Initially,
     only the top unit triangle is colored black, and all other unit triangles are colored white.

    A $\textbf{move}$ consists of choosing a black unit triangle and simultaneously flipping the colors of \textbf{all} unit triangles that share 
    an edge with it; that is, each such neighboring triangle changes from white to black or from black to white. Determine all positive integers
    $n$ for which it is possible, after a sequence of moves, for all unit triangles in the grid to be black at the same time.
\end{problem}

\begin{problem}
    Let $\sigma(n)$ denote the divisors sum of $n$, and let $d(n)$ denote the number of divisors of $n$. Prove that for any nonnegative 
    integer $k$, the equation $\sigma(n) = d(n)^k$ has finitely many solutions.
\end{problem}

\begin{problem}
    Let $\omega_1$ and $\omega_2$ be two circles intersecting at distinct points $X$ and $Y$. Let $\ell$ be their common tangent closer to $Y$, 
    intersecting $\omega_1$ at $A$ and $\omega_2$ at $B$. The tangent to $\omega_1$ at $Y$ meets $\omega_2$ again at $C$ with $C \neq Y$. 
    Let the line $XB$ meet $\omega_1$ again at $K$ with $K \neq X$. Let $D = BC \cap AY$, and let $E = YK \cap CX$.

    Prove that $D, Y, X, E$ are concyclic. (Four points are concyclic if there exists a circle containing all four points.)
\end{problem}

\begin{problem}
    For each positive integer $n$, let $\kappa(n)$ denote the unique squarefree positive integer such that$\frac{n}{\kappa(n)}$ is a perfect square.
    Let $G$ be a connected simple graph whose edges are labeled by positive integers $w_e$. Assume that for every cycle $C$ in $G$, the product
    $\prod_{e\in C} w_e$ is a perfect square. Prove that there exist squarefree positive integers $s_v$ assigned to the vertices $v\in V(G)$ 
    such that for every edge $uv\in E(G)$ there exists a positive integer $t_{uv}$ with $w_{uv}=\kappa(s_us_v)\,t_{uv}^2$.
\end{problem}

\begin{problem}
    For all positive reals $a,b,c$, prove the inequality
    $$\frac{a}{b+c}+\frac{b}{c+a}+\frac{c}{a+b} \; \ge\; \frac32+\frac{(a-b)^2+(b-c)^2+(c-a)^2}{2(a+b+c)^2}.$$
    Determine when equality holds.
\end{problem}

\newpage

\section{Solutions}

\begin{problem}
    Let $a_1,a_2,\dots,a_{2n+2}$ be positive real numbers. Prove that they can be permuted into numbers $u_0,v_0,u_1,v_1,\dots,u_n,v_n$ such that 
    the polynomial
    $$P(x)=u_0-v_0x+u_1x^2-v_1x^3+\dots+u_nx^{2n}-v_nx^{2n+1}$$
    has at least one real root, and every real root of $P$ is greater than or equal to $1$.
\end{problem}

\begin{proof}
    Obviously $P$ has one real root due to it having an odd degree. We then group the numbers into $n+1$ pairs like
    $(u_0, v_0)$, $(u_1, v_1)$, $\cdots (u_n, v_n)$ such that for each $0 \le k \le n$, $u_k \ge v_k$. We now observe the behavior of the
    following expression:

    \[
        u_kx^{2k} - v_kx^{2k+1} = x^{2k}((u_k-v_k)+v_k(1-x))
    \]

    Which, for $x<1$, is nonnegative. Moreover, consider the case when $x=0$, we get for $k=0$ that $u_0 - v_0(0) = u_0 > 0$ so
    at least one of these pairs is strictly positive for all $x<1$ and we conclude that

    \[
    P(x) > 0 \text{ for all } x < 1.
    \]

    FTSOC, if $r<1$ was a real root of $P$, we would have $P(r)=0$ and $P(r)>0$, contradiction. Therefore every real root $r$ of $P$ satisfies
    $r \ge 1$.
\end{proof}

\end{document}